\section{Methodology}
\label{sec:method}

We compare four retrieval configurations within a shared pipeline, isolating
the effect of post-retrieval decisions on answer faithfulness. The shared
infrastructure (Section~\ref{sec:shared}) ensures that differences in
outcomes are attributable solely to retrieval strategy.

\subsection{Shared pipeline}
\label{sec:shared}

Figure~\ref{fig:pipeline} shows the common architecture.
\input{../figures/pipeline}

\noindent All conditions share three components:
\begin{itemize}
  \item \textbf{Embedding:}
        \texttt{all-MiniLM-L6-v2}~\citep{reimers2019sbert}, a 384-dimensional
        sentence encoder (82M parameters), with vectors stored in ChromaDB
        using cosine similarity.
  \item \textbf{Generation:} Ollama-served
        Mistral-7B~\citep{jiang2023mistral} or
        Llama-3-8B~\citep{llama3}, both instruction-tuned, at temperature
        0.1 to minimize stochastic variation.
  \item \textbf{Faithfulness scoring:}
        \texttt{cross-encoder/nli-deberta-v3-base}~\citep{he2021deberta}
        (435M parameters). Each answer sentence is classified against the
        concatenated retrieved context; the mean entailment probability
        defines the faithfulness score in $[0,1]$. Responses below 0.50 are
        classified as hallucinated.
\end{itemize}

\subsection{Retrieval configurations}

\paragraph{Baseline (fixed chunking).}
Documents are split into fixed-length chunks of 256, 512, or 1024 tokens
with 10\% overlap using recursive character splitting. The top-$k$ passages
($k \in \{3, 5\}$) by cosine similarity are concatenated and passed to the
generator. This represents the standard RAG configuration against which all
interventions are measured.

\paragraph{HCPC-v1 (aggressive refinement).}
After retrieving 1024-token chunks, each passage is rescored with the
\texttt{ms-marco-MiniLM-L-6-v2} cross-encoder. If \emph{either} the cosine
similarity falls below 0.45 \emph{or} the cross-encoder logit falls below
$-0.20$, the chunk is split into 256-token sub-chunks and replaced by the
highest-scoring sub-span. This OR-gated policy is not a strawman: it
reflects the standard approach in production systems that apply
cross-encoder reranking with sub-chunk
replacement~\citep{xu2023recomp, nogueira2019passage} and is consistent
with the refinement logic in popular RAG frameworks that default to
aggressive post-retrieval filtering when any quality signal falls below a
threshold. We include HCPC-v1 as a representative instance of this
maximize-relevance philosophy, not as a deliberately weakened baseline.

\paragraph{HCPC-v2 (selective refinement).}
Our proposed method modifies the refinement policy in three principled ways,
each targeting a specific failure mode of HCPC-v1:

\begin{enumerate}
  \item \textbf{AND-gated detection.} Both the cosine similarity
        \emph{and} the cross-encoder score must fall below their respective
        thresholds ($\tau_{\text{sim}} = 0.45$, $\tau_{\text{ce}} = -0.20$)
        before a chunk is declared weak. This prevents refinement from
        triggering on passages that are borderline under one metric but
        adequate under the other, reducing false-positive interventions.

  \item \textbf{Rank protection.} The top-$N_{\text{prot}} = 2$ chunks by
        retrieval score are never refined, regardless of their cross-encoder
        scores. This protects the passages most likely to contain
        high-quality evidence from being fragmented.

  \item \textbf{Contiguity-preserving merge-back.} After splitting and
        rescoring, adjacent sub-chunks from the same parent passage are
        merged back into a longer span (up to 4096 characters $\approx$
        1024 tokens) when both pass the quality threshold. This restores
        the narrative continuity that splitting would otherwise destroy.
\end{enumerate}

Figure~\ref{fig:architecture} illustrates the full HCPC-v2 decision flow.
\input{../figures/architecture}

The design philosophy is deliberate conservatism. Each safeguard reduces the
frequency of refinement, which is precisely the mechanism through which
HCPC-v2 preserves coherence. On a strong retrieval set, the method
intervenes rarely; on a weak set, it intervenes more often but cannot fully
compensate for upstream encoder-corpus misalignment.

\paragraph{Cross-encoder re-ranking baseline.}
To separate refinement effects from scoring effects, we include a two-stage
reranker that over-fetches $k_{\text{fetch}} = 7$ candidates and returns the
top $k = 3$ by cross-encoder score, without any splitting or merging. This
baseline tests whether improved passage selection alone---without
refinement---improves faithfulness.

\subsection{The refinement paradox: why local optimization fails}

Figure~\ref{fig:failure_case} illustrates the mechanism behind the
refinement paradox with a concrete example.
\input{../figures/failure_case}

The failure occurs because retrieval-level optimization and generation-level
faithfulness operate at different granularities. A sub-chunk may score higher
against the query in isolation, but the information it replaces---topic
transitions, supporting details, contextual qualifiers---served as
connective tissue in the full context. The generator, which processes the
concatenated passage set as a sequence, loses the bridge sentences needed to
ground its reasoning and fills the gap from parametric memory.

This is not a pathological edge case. Under HCPC-v1, the OR gate triggers on
100\% of SQuAD queries despite the baseline retrieval being entirely adequate,
producing a systematic faithfulness degradation that standard retrieval
metrics would not detect.

\subsection{Evaluation framework}

We report six primary metrics:
\begin{enumerate}
  \item \textbf{NLI faithfulness} (mean entailment score, $[0,1]$).
  \item \textbf{Hallucination rate} (fraction of responses below 0.50).
  \item \textbf{Mean query-chunk similarity} (cosine between query and
        retrieved chunks).
  \item \textbf{CCS} (Equation~\ref{eq:ccs}; retrieval-time coherence).
  \item \textbf{Refinement rate} (fraction of queries triggering
        intervention).
  \item \textbf{Latency} (end-to-end generation time).
\end{enumerate}

Statistical analysis uses one-way ANOVA for chunk size, top-$k$, and prompt
effects, with effect sizes reported as $\eta^2$ and pairwise Cohen's $d$.
We use Spearman's $\rho$ for coherence-faithfulness correlations because
several retrieval metrics are not normally distributed.

\FloatBarrier
